\question
Утраченные отметки. Феникс Райт и Майя восстанавливают правило отметки улик $g(x,y,z)$ с независимыми двоичными признаками. Сохранился вектор $(??????01)$ в порядке $000,001,\ldots,111$; вопросительные знаки обозначают утраченные значения. В некоторой ДНФ функции имеется конъюнкт $\neg x$, а в некоторой КНФ --- дизъюнкт $\neg x\lor y$. Эти формы не обязательно совершенные.
\begin{enumerate}
\item Восстановите вектор и запишите номера всех минтермов и макстермов, начиная с нуля.
\item Определите принадлежность $g$ классам $T_0,T_1,S,M,L$. Положительные ответы обоснуйте; для отрицательных приведите контрпример или свойство полинома Жегалкина.
\end{enumerate}

---------------

Автор -- Олег Эмрих
